\documentclass[12pt]{article}
\usepackage[T1]{fontenc}
\usepackage[utf8]{inputenc}
\usepackage{lmodern}
\usepackage{textcomp}
\usepackage{enumitem}
\usepackage{geometry}
\geometry{margin=1in}
\begin{document}
\begin{center}
Answer Key - Sociology - Undergraduate Level Assessment 18 \
\small (Answers are ordered exactly as in the question paper.)
\end{center}

\section*{Multiple Choice}
\begin{enumerate}[label=\arabic*.]
\item \textbf{Answer:} B
\\ \textit{Options:} A) taking a stand on the issues neglected by feminism; B) studying society from the perspective of women; C) the recognition of difference and diversity in women's lives; D) a tendency to ignore the gendered nature of knowledge
\\ \textit{Note:} The term 'feminist standpoint' emphasizes the importance of understanding social structures and experiences from women's perspectives to highlight gender inequalities and challenge traditional viewpoints in sociology.
\item \textbf{Answer:} C
\\ \textit{Options:} A) the rates of violent crime were similar for men and women; B) women's sexual delinquency was more likely to be normalized than men's; C) women's criminal behaviour tended to reflect traditional gender roles; D) all of the above
\\ \textit{Note:} Smart's work highlighted that women's criminal behavior often arises from the constraints and expectations of traditional gender roles, illustrating how societal norms influence female offending patterns.
\item \textbf{Answer:} D
\\ \textit{Options:} A) access and exchange 'private' information about consumers; B) reduce prison overcrowding by the use of electronic tagging; C) monitor employees' activities at work; D) all of the above
\\ \textit{Note:} Technological forms of surveillance have made it easier to monitor, control, and analyze behavior in various contexts, thereby encompassing all potential implications that may arise from increased surveillance capabilities.
\item \textbf{Answer:} B
\\ \textit{Options:} A) white collar crime; B) organized crime; C) non-criminal deviance; D) global terrorism
\\ \textit{Note:} The Mafia exemplifies organized crime because it operates as a structured group involved in illegal activities for profit, often employing violence and corruption to maintain its power and influence within society.
\item \textbf{Answer:} B
\\ \textit{Options:} A) race is an objective way of categorizing people on biological grounds; B) the idea of race is socially constructed through powerful ideologies; C) race relations in Britain and America can be traced back to colonial times; D) people choose their racial identity and this becomes fixed
\\ \textit{Note:} The correct answer highlights that theories of racialized discourse emphasize that race is not a biological given but rather a concept shaped by social interactions and dominant ideologies that influence perceptions and definitions of racial identity.
\end{enumerate}

\section*{True / False}
\begin{enumerate}[label=\arabic*.]
\setcounter{enumi}{5}
\item \textbf{Answer:} False
\\ \textit{Options:} A) True; B) False
\\ \textit{Note:} Self-reported depression is actually found to be higher among women and those from lower socioeconomic backgrounds, as evidenced by the findings of Brown \& Harris (1978).
\item \textbf{Answer:} False
\\ \textit{Options:} A) True; B) False
\\ \textit{Note:} The statement is false because the nuclear family, consisting of two parents and their children, has existed in various forms throughout history, even prior to industrialization, alongside extended kin networks.
\item \textbf{Answer:} False
\\ \textit{Options:} A) True; B) False
\\ \textit{Note:} Berger (1963) actually employs the metaphor of a "theater" rather than a "circus" to describe social interactions and community dynamics, emphasizing the structured performances individuals engage in within social contexts.
\item \textbf{Answer:} False
\\ \textit{Options:} A) True; B) False
\\ \textit{Note:} Herberg's research indicated that ethnic minorities often maintain strong religious practices and community ties, which can lead to higher levels of engagement compared to mainstream populations.
\item \textbf{Answer:} False
\\ \textit{Options:} A) True; B) False
\\ \textit{Note:} Government 'white papers' are official documents intended to inform the public and policymakers about specific issues, and they are typically made accessible to a wider audience rather than being restricted to a select few.
\end{enumerate}

\section*{Long Form Response}
\begin{enumerate}[label=\arabic*.]
\setcounter{enumi}{10}
\item \textbf{Answer:} Social stratification refers to the hierarchical arrangement of individuals and groups in society based on various socio-economic factors. This structure is commonly depicted as layers, with the upper class at the top, followed by the middle class, and the working class at the bottom, each layer reflecting differing levels of wealth, power, and prestige. Factors influencing movement between these layers include education, income, social networks, and systemic barriers such as discrimination. For example, access to quality education can enable upward mobility, while economic downturns may trap individuals in lower strata. Overall, social stratification underscores the inequalities in resource distribution and opportunities available to members of society.
\\ \textit{Note:} correct because it accurately describes social stratification as a hierarchical system of layers based on socio-economic factors, identifies the classes within that structure, and outlines the various influences on social mobility, including education and systemic barriers.
\item \textbf{Answer:} Blauner's (1964) concept of alienation in the workplace, particularly in the car manufacturing industry, highlights how workers in assembly plants can feel disconnected from their labor and the products they create. This alienation arises from several factors, including the repetitive nature of assembly line work, strict supervision, and a lack of control over the production process. Workers often perform narrow tasks that do not allow for personal expression or creativity, leading to a sense of powerlessness and estrangement from their work. Additionally, the mechanization of tasks can contribute to feelings of insignificance, as workers may see themselves as mere cogs in a larger machine rather than valued contributors. Overall, these dynamics foster an environment where workers experience a profound disconnect from their labor, resulting in decreased job satisfaction and engagement.
\\ \textit{Note:} correct because it accurately identifies the key factors contributing to workplace alienation as described by Blauner, such as repetitive tasks, stringent supervision, and limited worker autonomy, specifically within the context of the car manufacturing industry.
\end{enumerate}

\vfill
\noindent\textit{End of Answer Key}
\end{document}